\chapter{La Struttura Sequenziale}
\label{chap:sequenza}

\begin{quote}
\itshape
\small
``Prima i muri, poi il tetto: in un algoritmo l'ordine dei passi non è un dettaglio, è la soluzione stessa.''
\end{quote}
\vspace{1em}

\noindent\textbf{Riferimento Lezione:} 28 Settembre 2026 -- Seconda Lezione di Laboratorio\\
\textbf{Competenze chiave:} Scomposizione di un problema nel modello Input--Elaborazione--Output, scrittura di sequenze lineari di istruzioni elementari, riconoscimento degli errori tipici di ordine e di ambiguità.

\vspace{1.5em}

% ------------------------------------------------------------------------------
\section{Ripartiamo dalla Definizione}
% ------------------------------------------------------------------------------

Nel capitolo precedente abbiamo definito l'algoritmo come una sequenza \textbf{finita} e \textbf{ordinata} di istruzioni elementari, chiare e \textbf{non ambigue}. Rileggiamo la definizione mettendo l'accento su tre parole:

\begin{itemize}
  \item \textbf{Sequenza ordinata:} l'ordine cronologico dei passi è vincolante. Scambiare due istruzioni altera il risultato o blocca l'esecuzione.
  \item \textbf{Istruzioni elementari:} passi atomici che l'esecutore (umano o macchina) sa compiere senza dubbi né interpretazioni soggettive.
  \item \textbf{Problema e soluzione:} l'algoritmo parte da una condizione iniziale (\textbf{Input}) e produce un risultato concreto (\textbf{Output}).
\end{itemize}

% ------------------------------------------------------------------------------
\section{Il Modello Fondamentale I-P-O}
% ------------------------------------------------------------------------------

Ogni elaborazione informatica si fonda su tre momenti, che d'ora in poi chiameremo \textbf{schema I-P-O} (\emph{Input -- Process -- Output}):

\vspace{0.5em}
\begin{center}
\begin{tcolorbox}[colback=white,colframe=PiazzaGray!40,width=0.88\textwidth,arc=2mm,halign=center]
\textbf{\color{BoxBlue}INPUT} $\longrightarrow$
\colorbox{PiazzaPrimary!10}{\textbf{\color{PiazzaPrimary}~ELABORAZIONE~}} $\longrightarrow$
\textbf{\color{BoxGreen}OUTPUT}
\end{tcolorbox}
\end{center}

\begin{enumerate}
  \item \textbf{Input (dati in ingresso):} le informazioni fornite al sistema, come valori digitati dall'utente, misure o prezzi.
  \item \textbf{Elaborazione (Process):} la serie ordinata di calcoli e trasformazioni eseguite dal calcolatore passo dopo passo.
  \item \textbf{Output (dati in uscita):} il risultato restituito, come un testo a video, un totale stampato o un disegno completato.
\end{enumerate}

\begin{consiglio}[Le Quattro Proprietà, in Breve]
Un algoritmo corretto è \textbf{finito} (termina sempre), \textbf{definito} (``aggiungi 5 grammi di zucchero'', non ``un po' di zucchero''), \textbf{generale} (risolve ogni problema della sua classe, non un solo caso con dati fissi) ed \textbf{efficiente}. In prima superiore la priorità è al 100\% sulla correttezza logica.
\end{consiglio}

% ------------------------------------------------------------------------------
\section{La Sequenza Pura}
% ------------------------------------------------------------------------------

Nelle prime settimane di laboratorio lavoriamo esclusivamente sulla \textbf{sequenza lineare}, la più semplice delle strutture di controllo.

\begin{concetto}[La Struttura Sequenziale]
In una \textbf{sequenza} le istruzioni vengono eseguite rigidamente \textbf{dall'alto verso il basso}, una dietro l'altra:
\begin{itemize}
  \item non ci sono bivi né ripetizioni: ogni riga viene eseguita \textbf{esattamente una volta};
  \item \textbf{l'ordine temporale è tassativo}: non si può usare o elaborare un dato prima di averlo letto o calcolato.
\end{itemize}
\end{concetto}

% ------------------------------------------------------------------------------
\section{Esempi Svolti}
% ------------------------------------------------------------------------------

\begin{esempio}[Costruire un Rifugio di Terra 3x3 in Minecraft]
\textbf{Input (inventario):} 33 blocchi di terra, 1 porta di legno, 1 torcia, 1 letto, 1 fornace, 1 banco da lavoro.\\
\textbf{Output:} rifugio $3 \times 3$ chiuso, illuminato e arredato per superare la notte.

\begin{enumerate}
  \item Posiziona il perimetro della base con 8 blocchi di terra, lasciando un vuoto per la porta.
  \item Innalza il secondo livello dei muri (8 blocchi).
  \item Innalza il terzo livello dei muri (8 blocchi).
  \item Copri il soffitto con 9 blocchi di terra: il tetto $3 \times 3$ è chiuso.
  \item Posiziona la porta di legno per sigillare l'ingresso.
  \item Attacca la torcia su una parete interna (la luce impedisce la comparsa dei mostri).
  \item Posiziona il letto in un angolo per impostare il punto di rinascita.
  \item Posiziona la fornace.
  \item Posiziona il banco da lavoro.
  \item Chiudi la porta e dormi: notte superata!
\end{enumerate}

\textbf{Perché l'ordine è legge?} Se provi a piazzare la torcia o il tetto prima di aver innalzato i muri, non hai un supporto solido su cui posizionarli. Se metti il letto prima di chiudere i muri, un mostro ti elimina prima dell'alba.
\end{esempio}

\begin{esempio}[Preparare un Toast Farcito alla Piastra]
\textbf{Input:} 2 fette di pane da toast, 2 fette di formaggio, 1 fetta di prosciutto cotto, piastra elettrica, 1 piatto.\\
\textbf{Output:} toast caldo e dorato con formaggio fuso, pronto sul piatto.

\begin{enumerate}
  \item Appoggia le due fette di pane sul tagliere pulito.
  \item Adagia la prima fetta di formaggio sulla fetta di pane inferiore.
  \item Aggiungi la fetta di prosciutto cotto sopra il formaggio.
  \item Adagia la seconda fetta di formaggio sopra il prosciutto.
  \item Chiudi il toast sovrapponendo la seconda fetta di pane.
  \item Accendi la piastra elettrica e attendi che sia calda.
  \item Posiziona il toast sulla piastra e chiudi il coperchio.
  \item Attendi 3 minuti per la tostatura.
  \item Apri la piastra e preleva il toast con una paletta.
  \item Adagia il toast sul piatto e servilo caldo.
\end{enumerate}

\textbf{Perché l'ordine è fondamentale?} Se metti il formaggio sulla piastra prima del pane, il formaggio brucia. Se provi a tostare il pane prima di accendere la piastra, resta crudo.
\end{esempio}

% ------------------------------------------------------------------------------
\newpage
\section{Esercizi di Laboratorio su Carta}
% ------------------------------------------------------------------------------

\begin{esercizio}[Esercizio 1: La Media di 3 Voti]
Uno studente ha ottenuto 3 voti in informatica: \textbf{7}, \textbf{8} e \textbf{6}.\\
Scrivi sul quaderno, in italiano chiaro, la sequenza lineare che calcola e restituisce la media esatta. Indica prima l'\textbf{Input}, poi i passi dell'elaborazione e infine l'\textbf{Output}.
\end{esercizio}

\begin{soluzione}[Esercizio 1: La Media di 3 Voti]
\textbf{Input:} i tre voti (7, 8 e 6).
\begin{enumerate}
  \item Leggi il primo voto, il secondo voto e il terzo voto.
  \item Calcola la somma dei tre voti: $7 + 8 + 6 = 21$.
  \item Dividi la somma per il numero dei voti: $21 \div 3 = 7$.
  \item Mostra la media ottenuta.
\end{enumerate}
\textbf{Output:} la media, cioè 7.

\smallskip
\noindent Per rispettare la \textbf{generalità} la sequenza deve funzionare con qualsiasi terna di voti: si scrive ``somma dei tre voti diviso 3'', non ``21 diviso 3''.
\end{soluzione}

\begin{esercizio}[Esercizio 2: Il Robot Disegnatore]
Un robot su una lavagna digitale comprende solo due istruzioni:
\begin{itemize}
  \item \textbf{``Avanza di 100 passi tracciando una linea dritta''};
  \item \textbf{``Gira a destra di 90 gradi''}.
\end{itemize}
Scrivi la sequenza esatta di 8 istruzioni che fa disegnare al robot un quadrato chiuso e lo fa tornare orientato come in partenza.
\end{esercizio}

\begin{soluzione}[Esercizio 2: Il Robot Disegnatore]
\begin{enumerate}
  \item Avanza di 100 passi tracciando una linea dritta.
  \item Gira a destra di 90 gradi.
  \item Avanza di 100 passi tracciando una linea dritta.
  \item Gira a destra di 90 gradi.
  \item Avanza di 100 passi tracciando una linea dritta.
  \item Gira a destra di 90 gradi.
  \item Avanza di 100 passi tracciando una linea dritta.
  \item Gira a destra di 90 gradi.
\end{enumerate}
Le istruzioni 1--4 tracciano due lati; le 5--8 chiudono il quadrato. L'ultima rotazione (passo 8) non disegna nulla, ma riporta il robot nella direzione di partenza: senza di essa il quadrato sarebbe completo, ma la consegna no.
\end{soluzione}

% ------------------------------------------------------------------------------
\section{I 4 Errori Tipici della Sequenza}
% ------------------------------------------------------------------------------

\begin{attenzione}[Gli Errori da Evitare]
\begin{enumerate}
  \item \textbf{Dare per scontato l'ovvio:} saltare un passaggio intermedio che per noi è banale ma per la macchina è indispensabile.
  \item \textbf{Passi invertiti:} calcolare il totale prima di aver letto i dati in ingresso.
  \item \textbf{Istruzioni ambigue:} usare termini vaghi come ``un po''' invece di numeri e misure esatte.
  \item \textbf{Soluzioni cablate (\emph{hardcoded}):} scrivere il risultato già calcolato a mente anziché il procedimento generale.
\end{enumerate}
\end{attenzione}
